% ============================================================
%  CUADERNILLO DE ESTUDIO — CLASE 2
%  Seguridad de la Información y de Sistemas
%  Lic. en Gestión de Negocios Digitales — UCA
%  Compilar con: pdflatex  (dos pasadas)
% ============================================================

\documentclass[12pt,a4paper]{article}

% ---------- Paquetes ----------
\usepackage[utf8]{inputenc}
\usepackage[spanish]{babel}
\usepackage[T1]{fontenc}
\usepackage{lmodern}
\usepackage[margin=2.2cm]{geometry}
\usepackage{amssymb}
\usepackage{enumitem}
\usepackage{booktabs}
\usepackage{array}
\usepackage{tabularx}
\usepackage{xcolor}
\usepackage{tcolorbox}
\usepackage{graphicx}
\usepackage{fancyhdr}
\usepackage{titlesec}
\usepackage{hyperref}
\usepackage{multicol}
\usepackage{pifont}
\usepackage{wasysym}
\usepackage{tikz}
\usetikzlibrary{positioning,arrows.meta}

% ---------- Colores ----------
\definecolor{azulUCA}{HTML}{003366}
\definecolor{doradoUCA}{HTML}{B8860B}
\definecolor{grisClaro}{HTML}{F2F2F2}
\definecolor{rojoAlerta}{HTML}{C0392B}
\definecolor{verdeOK}{HTML}{27AE60}
\definecolor{azulCielo}{HTML}{2980B9}

% ---------- tcolorbox estilos ----------
\tcbuselibrary{skins,breakable}

\newtcolorbox{cajaTitulo}[1][]{
  colback=azulUCA, colframe=azulUCA, coltext=white,
  fonttitle=\Large\bfseries, arc=4mm, #1
}

\newtcolorbox{cajaDefinicion}{
  colback=azulCielo!10, colframe=azulCielo,
  fonttitle=\bfseries, title=Definición, arc=3mm, breakable
}

\newtcolorbox{cajaEjemplo}{
  colback=verdeOK!10, colframe=verdeOK,
  fonttitle=\bfseries, title=Ejemplo, arc=3mm, breakable
}

\newtcolorbox{cajaReflexion}{
  colback=doradoUCA!10, colframe=doradoUCA,
  fonttitle=\bfseries, title=Reflexión, arc=3mm, breakable
}

\newtcolorbox{cajaActividad}{
  colback=grisClaro, colframe=azulUCA,
  fonttitle=\bfseries\color{azulUCA}, title={\ding{45}\; Actividad}, arc=3mm, breakable
}

\newtcolorbox{cajaNota}{
  colback=rojoAlerta!8, colframe=rojoAlerta,
  fonttitle=\bfseries\color{rojoAlerta}, title=Nota, arc=3mm, breakable
}

% ---------- Header / Footer ----------
\pagestyle{fancy}
\fancyhf{}
\fancyhead[L]{\small\textcolor{azulUCA}{Seguridad de la Información y de Sistemas}}
\fancyhead[R]{\small\textcolor{azulUCA}{Cuadernillo — Clase 2}}
\fancyfoot[C]{\thepage}
\fancyfoot[L]{\small\textcolor{gray}{UCA — 2.º cuatr.}}
\renewcommand{\headrulewidth}{0.4pt}
\renewcommand{\footrulewidth}{0.4pt}

% ---------- Títulos de sección ----------
\titleformat{\section}
  {\color{azulUCA}\Large\bfseries}
  {\thesection.}{0.6em}{}
\titleformat{\subsection}
  {\color{azulCielo}\large\bfseries}
  {\thesubsection}{0.6em}{}

% ---------- Enlaces ----------
\hypersetup{colorlinks=true, linkcolor=azulUCA, urlcolor=azulCielo}

% ---------- Checkboxes ----------
\newcommand{\checkbox}{\ding{111}\;}
\newcommand{\checkboxVacio}{$\square$\;}

% ============================================================
\begin{document}

% ============================================================
%  PORTADA
% ============================================================
\begin{titlepage}
\centering
\vspace*{1cm}

% Escudo / Título principal
\begin{cajaTitulo}[width=\textwidth, halign=center]
  SEGURIDAD DE LA INFORMACIÓN Y DE SISTEMAS\\[4pt]
  {\large Licenciatura en Gestión de Negocios Digitales — UCA}\\[2pt]
  {\normalsize 2.º Cuatrimestre}
\end{cajaTitulo}

\vspace{1.5cm}

{\Huge\bfseries\color{azulUCA} CUADERNILLO DE ESTUDIO}\\[8pt]
{\LARGE\color{doradoUCA} CLASE 2}\\[12pt]
{\Large\itshape ``Fundamentos de Seguridad (CID)\\[2pt] y Debrief de la Huella Digital''}\\[1.5cm]

% Datos
\begin{tabular}{rl}
\textbf{Profesor:}   & Mg.\ Ing.\ Rodrigo Atilio Elgueta \\[4pt]
%\textbf{Fecha:}      & Jueves 13 de agosto de 2026 \\[4pt]
\textbf{Unidad:}     & Unidad 1 \\[4pt]
\end{tabular}

\vfill

% Datos del estudiante
\begin{tcolorbox}[colback=grisClaro, colframe=azulUCA, arc=3mm, width=0.85\textwidth]
  \centering
  {\bfseries Datos del/la estudiante}\\[8pt]
  \begin{tabular}{p{5cm}l p{5cm}l}
    Nombre y apellido: & \underline{\hspace{4.8cm}} &
    Grupo / N.º: & \underline{\hspace{4.8cm}} \\[10pt]
    Correo: & \underline{\hspace{4.8cm}} &
    Fecha: & \underline{\hspace{4.8cm}} \\
  \end{tabular}
\end{tcolorbox}

\vspace{1cm}
{\small\textcolor{gray}{Material de uso exclusivo para la cursada. Prohibida su distribución externa.}}
\end{titlepage}

% ============================================================
%  ÍNDICE (opcional)
% ============================================================
\tableofcontents
\newpage

% ============================================================
%  SECCIÓN 1 — DEBRIEF
% ============================================================
\section{¿Qué vivimos la semana pasada? (Debrief)}

La Clase~1 ingresaste desde tu celular a una página web que parecía inofensiva.
Sin que lo notaras de forma clara, esa página \textbf{accedió a la cámara de tu
dispositivo} y transmitió la imagen a un panel de administración controlado por
el docente.

\begin{cajaReflexion}
No hubo un ``hackeo'' sofisticado. Hubo un diseño que aprovechó que no leíste,
no notaste o aceptaste un permiso sin pensar. Eso, en un negocio digital real,
es una \textbf{brecha de seguridad}.
\end{cajaReflexion}

\begin{cajaActividad}[title={\ding{45}\; Actividad individual — Completá durante la clase}]
Respondé con tus palabras en cada recuadro:

\vspace{6pt}
\begin{tabularx}{\textwidth}{|>{\bfseries}p{0.32\textwidth}|X|}
\hline
\textbf{① ¿Qué dato o permiso cediste sin advertirlo plenamente?} & \\[2.5cm]
\hline
\textbf{② ¿Qué mecanismo usó la página para obtenerlo?} \newline
\footnotesize{(Pista: permisos del navegador, botones de ``aceptar'', redirecciones\ldots)} & \\[2.5cm]
\hline
\textbf{③ ¿Qué medida concreta habría prevenido o mitigado esto?} & \\[2.5cm]
\hline
\textbf{④ ¿A cuál concepto de hoy lo asociás?} \newline
\footnotesize{(activo, amenaza, vulnerabilidad, impacto, riesgo)} & \\[2.5cm]
\hline
\end{tabularx}
\end{cajaActividad}

\newpage
% ============================================================
%  SECCIÓN 2 — LOS 5 CONCEPTOS
% ============================================================
\section{Los 5 conceptos fundamentales}

Antes de hablar de ``seguridad informática'' hay que dominar cinco palabras.
Todo el curso ---y toda la práctica profesional--- se apoya en ellas.

% --- 2.1 Activo ---
\subsection{Activo de información}
\begin{cajaDefinicion}
Cualquier elemento que tenga valor para la organización y deba protegerse.
\end{cajaDefinicion}

No es solo ``el servidor''. Son activos:
\begin{itemize}[leftmargin=2em]
  \item Los \textbf{datos} de clientes (nombres, tarjetas, historial de compras).
  \item El \textbf{software} propio (código de una app, algoritmos).
  \item El \textbf{hardware} (servidores, notebooks, celulares corporativos).
  \item El \textbf{conocimiento} (manuales de proceso, estrategias comerciales).
  \item La \textbf{reputación} de la marca.
  \item Las \textbf{personas} (sus credenciales, su acceso).
\end{itemize}

\begin{cajaEjemplo}
\textbf{Para un e-commerce:} la base de datos con 50\,000 clientes, el código de
la pasarela de pago, el servidor donde corre la tienda, la cuenta de Instagram
con 200\,000 seguidores.
\end{cajaEjemplo}

% --- 2.2 Amenaza ---
\subsection{Amenaza}
\begin{cajaDefinicion}
Un evento o actor potencial que puede causar daño a un activo.
\end{cajaDefinicion}

La amenaza \textbf{no es el daño en sí}, es la \emph{posibilidad} de que algo malo pase.
\begin{itemize}[leftmargin=2em]
  \item Un atacante que envía correos de \textbf{phishing}.
  \item Un \textbf{ransomware} que cifra los discos.
  \item Un empleado que hace clic donde no debe (\textbf{ingeniería social}).
  \item Un desastre natural que inunda el data center.
  \item Un competidor que hace \textbf{scraping} agresivo de tu catálogo.
\end{itemize}

% --- 2.3 Vulnerabilidad ---
\subsection{Vulnerabilidad}
\begin{cajaDefinicion}
Una debilidad en un activo, proceso o control que puede ser explotada por una amenaza.
\end{cajaDefinicion}

La vulnerabilidad \textbf{está en tu lado}, no en el atacante.
\begin{itemize}[leftmargin=2em]
  \item Un software sin parches de seguridad.
  \item Una contraseña débil como \texttt{123456}.
  \item \textbf{No leer los permisos} que pide una página o app (¡lo que pasó en la Clase~1!).
  \item Un empleado que no recibió capacitación.
  \item Un bucket de almacenamiento en la nube configurado como ``público''.
\end{itemize}

% --- 2.4 Impacto ---
\subsection{Impacto}
\begin{cajaDefinicion}
La consecuencia (económica, legal, reputacional, operativa) que se materializa
cuando una amenaza explota una vulnerabilidad.
\end{cajaDefinicion}

\begin{table}[h!]
\centering
\begin{tabularx}{\textwidth}{lX}
\toprule
\textbf{Tipo} & \textbf{Ejemplo} \\
\midrule
Económico & Multa de la AAIP por filtración; pérdida de ventas durante 48\,h de caída. \\
Legal     & Demanda colectiva; sanción por incumplimiento de Ley 25.326. \\
Reputacional & Tendencia negativa en redes; pérdida de confianza. \\
Operativo & No se pueden procesar pagos durante un fin de semana largo. \\
\bottomrule
\end{tabularx}
\caption{Tipos de impacto.}
\end{table}

% --- 2.5 Riesgo ---
\subsection{Riesgo}
\begin{cajaDefinicion}
La combinación de la \textbf{probabilidad} de que una amenaza explote una
vulnerabilidad y el \textbf{impacto} que eso generaría.
\end{cajaDefinicion}

\begin{center}
\fbox{\parbox{0.75\textwidth}{\centering
\Large $\displaystyle \text{RIESGO} \;\approx\; \text{PROBABILIDAD} \;\times\; \text{IMPACTO}$
}}
\end{center}

\begin{cajaEjemplo}[title={Ejemplo integrado}]
Para un e-commerce argentino chico:
\begin{itemize}[leftmargin=2em]
  \item \textbf{Activo:} base de datos de 5\,000 clientes con datos de tarjeta.
  \item \textbf{Amenaza:} ataque de phishing a un empleado de atención al cliente.
  \item \textbf{Vulnerabilidad:} el empleado no fue capacitado y usa la misma
        contraseña en el mail corporativo y en sitios personales.
  \item \textbf{Impacto:} robo de datos de tarjetas $\rightarrow$ multa + pérdida
        de clientes + nota en medios.
  \item \textbf{Riesgo:} ALTO, porque la probabilidad es media-alta y el impacto
        es severo.
\end{itemize}
\end{cajaEjemplo}

% --- Esquema visual ---
\subsection*{Esquema resumen}
\begin{center}
\begin{tikzpicture}[
  node distance=1.8cm,
  box/.style={rectangle, draw=azulUCA, fill=azulUCA!10,
              text width=3cm, align=center, rounded corners=3pt,
              font=\small\bfseries},
  arr/.style={-{Stealth[length=3mm]}, thick, azulUCA}
]
  \node[box] (amen) {AMENAZA};
  \node[box, right=2.2cm of amen] (vuln) {VULNERABILIDAD};
  \node[box, right=2.2cm of vuln] (act)  {ACTIVO};
  \node[box, below=1.4cm of act]  (imp)  {IMPACTO};
  \node[box, below=1.4cm of imp,
        fill=doradoUCA!15, draw=doradoUCA] (ries) {RIESGO};

  \draw[arr] (amen) -- node[above,font=\scriptsize]{explota} (vuln);
  \draw[arr] (vuln) -- node[above,font=\scriptsize]{afecta} (act);
  \draw[arr] (act)  -- (imp);
  \draw[arr] (imp)  -- node[right,font=\scriptsize]{Prob.\ $\times$ Impacto} (ries);
\end{tikzpicture}
\end{center}

\newpage
% ============================================================
%  SECCIÓN 3 — PILARES CID
% ============================================================
\section{Los tres pilares CID}

La seguridad de la información se sostiene sobre tres propiedades.
Si alguna se rompe, hay un \textbf{incidente}.

% --- 3.1 Confidencialidad ---
\subsection{Confidencialidad \;\textrm{\Large\lock}}
\begin{cajaDefinicion}
La información solo es accesible para quienes están \textbf{autorizados}.
\end{cajaDefinicion}

\begin{table}[h!]
\centering
\begin{tabularx}{\textwidth}{XX}
\toprule
\textbf{Se cumple cuando\ldots} & \textbf{Se viola cuando\ldots} \\
\midrule
Solo RR.\,HH.\ ve los sueldos. &
Un empleado reenvía la nómina por error a todo el equipo. \\[4pt]
La conexión usa HTTPS. &
La demo de Clase 1: la cámara se activó sin consentimiento. \\[4pt]
Cifrado extremo a extremo. &
Un bucket de S3 queda ``público'' y cualquiera descarga datos. \\
\bottomrule
\end{tabularx}
\end{table}

% --- 3.2 Integridad ---
\subsection{Integridad \;\textrm{\Large\checkmark}}
\begin{cajaDefinicion}
La información es \textbf{exacta y completa}, y no fue modificada sin autorización.
\end{cajaDefinicion}

\begin{table}[h!]
\centering
\begin{tabularx}{\textwidth}{XX}
\toprule
\textbf{Se cumple cuando\ldots} & \textbf{Se viola cuando\ldots} \\
\midrule
Un hash SHA-256 coincide con el original. &
Un atacante modifica precios en la base de datos. \\[4pt]
Registros contables sin alteraciones. &
Un ransomware cifra y corrompe archivos. \\[4pt]
Un correo llega sin modificaciones. &
Alguien edita una factura PDF antes de reenviarla. \\
\bottomrule
\end{tabularx}
\end{table}

% --- 3.3 Disponibilidad ---
\subsection{Disponibilidad \;\textrm{\Large\clock}}
\begin{cajaDefinicion}
La información y los sistemas están \textbf{accesibles cuando se los necesita}.
\end{cajaDefinicion}

\begin{table}[h!]
\centering
\begin{tabularx}{\textwidth}{XX}
\toprule
\textbf{Se cumple cuando\ldots} & \textbf{Se viola cuando\ldots} \\
\midrule
El e-commerce funciona un Black Friday. &
Un DDoS deja la tienda offline 12 horas. \\[4pt]
Hay backups probados y plan de recuperación. &
Se borra la base y no hay backup hace 3 meses. \\[4pt]
Pagos procesados 24/7. &
Corte de energía sin UPS ni generador. \\
\bottomrule
\end{tabularx}
\end{table}

% --- Regla mnemotécnica ---
\begin{tcolorbox}[colback=doradoUCA!10, colframe=doradoUCA, title=Regla mnemotécnica, arc=3mm]
\begin{center}
\textbf{C}onfidencialidad $\rightarrow$ ¿Quién lo ve?\\[4pt]
\textbf{I}ntegridad $\rightarrow$ ¿Está intacto?\\[4pt]
\textbf{D}isponibilidad $\rightarrow$ ¿Puedo usarlo cuando lo necesito?
\end{center}
\end{tcolorbox}

% --- Ejercicio ---
\begin{cajaActividad}[title={\ding{45}\; Ejercicio rápido — Marcá con C, I o D}]
\vspace{4pt}
\renewcommand{\arraystretch}{1.8}
\begin{tabularx}{\textwidth}{|c|X|c|}
\hline
\textbf{\#} & \textbf{Escenario} & \textbf{Pilar} \\
\hline
1 & Un ex empleado publica en redes la lista de clientes. & \checkboxVacio \\
\hline
2 & Un malware modifica cuentas bancarias destino en una fintech. & \checkboxVacio \\
\hline
3 & Un DDoS deja caída la app de delivery toda la hora del almuerzo. & \checkboxVacio \\
\hline
4 & Base de datos cifrada pero legible por cualquier empleado. & \checkboxVacio \\
\hline
5 & El servidor de backups se corrompe; 6 meses ilegibles. & \checkboxVacio \\
\hline
6 & Interceptación de HTTP sin cifrar; se leen datos de tarjeta. & \checkboxVacio \\
\hline
\end{tabularx}
\end{cajaActividad}

\newpage
% ============================================================
%  SECCIÓN 4 — CONEXIÓN CON LA DEMO
% ============================================================
\section{Conectar la demo con la teoría}

Volvamos a lo que pasó la Clase~1 y clasifiquémoslo:

\begin{table}[h!]
\centering
\renewcommand{\arraystretch}{1.5}
\begin{tabularx}{\textwidth}{|>{\bfseries}p{3.2cm}|X|}
\hline
Activo       & Tu imagen personal, tu cámara, tu privacidad. \\
\hline
Amenaza      & El sitio web diseñado para acceder a la cámara sin evidencia clara. \\
\hline
Vulnerabilidad & Aceptar permisos sin leer; no verificar la URL; confianza ciega en un enlace. \\
\hline
Impacto      & Tu foto terminó en un panel que no controlás. En escala real: chantaje, suplantación, exposición pública. \\
\hline
Riesgo       & Alto si la página fuera maliciosa y el activo fueran datos sensibles. \\
\hline
Pilar CID violado & \textbf{Confidencialidad} (se accedió a tu imagen sin autorización consciente). \\
\hline
\end{tabularx}
\end{table}

\begin{cajaNota}[title={Lección para negocios digitales}]
Si una página ``inocente'' puede hacer esto, pensá qué puede hacer un
competidor, un ciberdelincuente o un empleado interno malintencionado con
acceso a los activos de tu empresa. La seguridad no es un gasto: es un
\textbf{habilitador de confianza} (concepto clave del programa).
\end{cajaNota}

% ============================================================
%  SECCIÓN 5 — GLOSARIO
% ============================================================
\section{Glosario mínimo de la Clase 2}

\begin{table}[h!]
\centering
\renewcommand{\arraystretch}{1.4}
\begin{tabularx}{\textwidth}{>{\bfseries}p{4cm}X}
\toprule
Término & Definición en una línea \\
\midrule
Activo de información & Todo lo que tiene valor y debe protegerse. \\
Amenaza               & Evento o actor que \emph{puede} causar daño. \\
Vulnerabilidad        & Debilidad que la amenaza puede explotar. \\
Impacto               & Consecuencia del daño materializado. \\
Riesgo                & Probabilidad $\times$ Impacto. \\
Confidencialidad      & Solo accede quien está autorizado. \\
Integridad            & La información no fue alterada indebidamente. \\
Disponibilidad        & La información está accesible cuando se necesita. \\
Huella digital        & Rastro de datos que dejamos al usar tecnología. \\
Permiso (app/navegador) & Autorización para acceder a un recurso del dispositivo. \\
\bottomrule
\end{tabularx}
\end{table}

\newpage
% ============================================================
%  SECCIÓN 6 — CUESTIONARIO 1
% ============================================================
\section{Cuestionario 1 de autoevaluación}

\begin{cajaNota}[title={Instrucciones}]
Este cuestionario se realiza al cierre de la clase (Kahoot / Forms / papel).
\textbf{No es eliminatorio.} Usalo para medir cuánto entendiste antes de seguir.
Las respuestas correctas están marcadas con \ding{51} al final de la sección.
\end{cajaNota}

\vspace{6pt}

% Pregunta 1
\textbf{1.} ¿Qué significa la sigla CID en seguridad de la información?
\begin{enumerate}[label=\alph*), leftmargin=2.5em]
  \item Control, Inspección, Detección
  \item Confidencialidad, Integridad, Disponibilidad
  \item Ciberseguridad, Información, Datos
  \item Confianza, Identidad, Datos
\end{enumerate}

% Pregunta 2
\textbf{2.} Un activo de información es:
\begin{enumerate}[label=\alph*), leftmargin=2.5em]
  \item Solo el hardware de una empresa
  \item Cualquier elemento que tenga valor para la organización y deba protegerse
  \item Únicamente los datos personales de los clientes
  \item El software antivirus instalado
\end{enumerate}

% Pregunta 3
\textbf{3.} La diferencia entre amenaza y vulnerabilidad es:
\begin{enumerate}[label=\alph*), leftmargin=2.5em]
  \item Son sinónimos
  \item La amenaza es un evento potencial dañino; la vulnerabilidad es una debilidad que puede ser explotada
  \item La vulnerabilidad es externa y la amenaza es interna
  \item No hay diferencia relevante para el negocio
\end{enumerate}

% Pregunta 4
\textbf{4.} En la demo de la Clase~1, ¿qué permiso del dispositivo se solicitó
sin que el usuario lo notara conscientemente?
\begin{enumerate}[label=\alph*), leftmargin=2.5em]
  \item Ubicación GPS
  \item Cámara
  \item Contactos
  \item Micrófono
\end{enumerate}

% Pregunta 5
\textbf{5.} La ``huella digital'' se refiere a:
\begin{enumerate}[label=\alph*), leftmargin=2.5em]
  \item Solo las huellas dactilares usadas en biometría
  \item El rastro de datos e información que dejamos al usar tecnología
  \item Un tipo de malware
  \item La firma digital de un documento
\end{enumerate}

% Pregunta 6
\textbf{6.} ¿Cuál describe mejor un \textbf{riesgo} (y no solo una amenaza)
para un e-commerce?
\begin{enumerate}[label=\alph*), leftmargin=2.5em]
  \item Que exista phishing en internet en general
  \item La probabilidad de que un ataque de phishing exitoso cause pérdida de datos de tarjetas y sanciones regulatorias
  \item Tener una vulnerabilidad de software sin explotar
  \item Ninguna de las anteriores
\end{enumerate}

% Pregunta 7
\textbf{7.} La seguridad de la información en un negocio digital debe entenderse
principalmente como:
\begin{enumerate}[label=\alph*), leftmargin=2.5em]
  \item Un gasto que hay que minimizar
  \item Un habilitador de confianza y una ventaja competitiva
  \item Un tema exclusivo del área de sistemas
  \item Una obligación solo para empresas grandes
\end{enumerate}

% Pregunta 8
\textbf{8.} Verdadero o Falso: la disponibilidad se refiere a que la información
esté accesible solo para quien está autorizado.

\vspace{4pt}
\checkboxVacio Verdadero \qquad \checkboxVacio Falso

% --- Clave de respuestas ---
\vspace{12pt}
\begin{tcolorbox}[colback=verdeOK!10, colframe=verdeOK,
                  title={\ding{51}\; Clave de respuestas}, arc=3mm]
\begin{enumerate}[leftmargin=2em]
  \item \textbf{b)} Confidencialidad, Integridad, Disponibilidad
  \item \textbf{b)} Cualquier elemento que tenga valor\ldots
  \item \textbf{b)} La amenaza es un evento potencial dañino\ldots
  \item \textbf{b)} Cámara
  \item \textbf{b)} El rastro de datos e información\ldots
  \item \textbf{b)} La probabilidad de que un ataque de phishing\ldots
  \item \textbf{b)} Un habilitador de confianza\ldots
  \item \textbf{Falso.} Eso describe a la confidencialidad; la disponibilidad es
        que la información esté accesible \emph{cuando se la necesita}.
\end{enumerate}
\end{tcolorbox}

\newpage
% ============================================================
%  SECCIÓN 7 — RESPUESTAS DEL EJERCICIO CID
% ============================================================
\section{Respuestas del ejercicio de la Sección 3}

\begin{table}[h!]
\centering
\renewcommand{\arraystretch}{1.5}
\begin{tabularx}{\textwidth}{|c|X|>{\centering\arraybackslash}p{3.5cm}|}
\hline
\textbf{\#} & \textbf{Escenario} & \textbf{Pilar afectado} \\
\hline
1 & Ex empleado publica lista de clientes & \textbf{C} \\
\hline
2 & Malware modifica cuentas bancarias destino & \textbf{I} \\
\hline
3 & DDoS deja caída la app de delivery & \textbf{D} \\
\hline
4 & Base cifrada pero legible por cualquier empleado & \textbf{C} \\
\hline
5 & Backups de 6 meses corruptos & \textbf{D} e \textbf{I} \\
\hline
6 & Interceptación de HTTP sin cifrar & \textbf{C} \\
\hline
\end{tabularx}
\end{table}

% ============================================================
%  SECCIÓN 8 — RECURSOS
% ============================================================
\section{Para seguir profundizando (optativo)}

\begin{table}[h!]
\centering
\renewcommand{\arraystretch}{1.4}
\begin{tabularx}{\textwidth}{>{\bfseries}p{5.5cm}X}
\toprule
Recurso & Qué vas a encontrar \\
\midrule
ISO/IEC 27000 (glosario) & Definiciones formales de activo, amenaza, vulnerabilidad, riesgo. \\
\texttt{haveibeenpwned.com} & Verificá si tu mail apareció en alguna brecha de datos pública. \\
INCIBE --- Guía de ciberseguridad para pymes & Lenguaje simple, ejemplos prácticos. \\
Video: ``HTTP vs HTTPS'' & Para entender por qué el cifrado protege la confidencialidad. \\
\bottomrule
\end{tabularx}
\end{table}

% ============================================================
%  SECCIÓN 9 — NOTAS
% ============================================================
\section{Espacio para notas personales}

\begin{tcolorbox}[colback=grisClaro, colframe=gray!50, arc=2mm, height=8cm]
\textcolor{gray}{\itshape Usá este espacio durante la clase para anotar dudas,
ejemplos útiles o preguntas para la próxima clase.}
\end{tcolorbox}

% ============================================================
%  SECCIÓN 10 — PRÓXIMA CLASE
% ============================================================
\section{¿Qué sigue? (Puente a la Clase 3)}

La próxima clase vamos a:
\begin{itemize}[leftmargin=2em]
  \item Ver el \textbf{panorama de amenazas} más comunes para negocios digitales
        (phishing, ransomware, ingeniería social).
  \item Probar \textbf{herramientas gratuitas reales}: Have I Been Pwned,
        VirusTotal, un gestor de contraseñas.
  \item \textbf{Conformar los grupos de 4-5} que van a trabajar todo el
        cuatrimestre, incluida la consigna del \textbf{Trabajo Práctico Integrador}.
\end{itemize}

\begin{center}
\begin{tcolorbox}[colback=rojoAlerta!8, colframe=rojoAlerta,
                  width=0.8\textwidth, halign=center, arc=3mm]
  \Large\ding{42}\; \textbf{¡Traé el celular con batería!}\\[4pt]
  \normalsize Vamos a usar herramientas en vivo.
\end{tcolorbox}
\end{center}

\vfill
\begin{center}
\small\textcolor{gray}{%
Cuadernillo elaborado para la cátedra de Seguridad de la Información y de
Sistemas --- Lic.\ en Gestión de Negocios Digitales, UCA, 2.º cuatrimestre.\\
Prohibida su distribución fuera del ámbito de la cursada sin autorización del docente.}
\end{center}

\end{document}